If $\mathfrak{U}$ is a subgroup of $\mathfrak{G}$ and if the factor group $\mathfrak{G}/\mathfrak{U}$ has $r_0$ basis elements, then among the $h$ characters of $\mathfrak{G}$ there are $r_0$ characters $\chi_i$ with order $h_i$, a power of a prime, $(i = 1, 2, \ldots, r_0)$ such that the $r_0$ conditions

$$\chi_i(A) = 1 \quad (i = 1, 2, \ldots, r_0)$$

are satisfied for all elements $A$ of $\mathfrak{U}$ and only for elements $A$ of $\mathfrak{U

A system of finitely or infinitely many elements of $\mathfrak{G}$: $A_i (i = 1, 2, \ldots)$, $T_k (k = 1, 2, \ldots)$ ($A_i$ of infinite order, $T_k$ of finite order), is called a basis for $\mathfrak{G}$ if each element of $\mathfrak{G}$ can be represented in the form

$$A_1^{x_1} A_2^{x_2} \cdots T_1^{y_1} T_2^{y_2} \cdots = C,$$

where

(1) the exponents $x_i$ and $y_k$ are integers and only finitely many are $\neq 0$,

(2) the exponents $x_i$ are determined uniquely and the exponents $y_k$ are determined uniquely mod $h_k$ by $C$.

Obviously any finite set of elements of a basis must be independent.

The requirement that the $h_k$ are powers of a prime will not be imposed here for the sake of simplicity.

A basis is called finite if it consists of finitely many elements.
